\Article{28}
Combining formulas I....IV gives the variation of the expression $W$
\[
\begin{aligned}
\delta W={}&\int\mathrm{d}U.\delta e.\cos(4,5).[z+\alpha\alpha(\frac{1}{R}+\frac{1}{R'})]\\
&-\int\mathrm{d}P.\delta e.\cos(5,8).(\alpha\alpha\cos i-\alpha\alpha+2\beta\beta)
\end{aligned}
\]
where the first integral must extend over all elements $\mathrm{d}U$ of the free part of the surface of the region $s$, or of the free parts if there happen to be several separate ones. The second integral must extend over all elements $\mathrm{d}P$ of the line or lines that separate that free part or those free parts from the remaining parts adjoining the region $S$.

Now the value of $W$ must be a minimum in the state of equilibrium and therefore cannot undergo a negative change for any infi%
\SourceLeaf{72}nitely small change in the shape of the fluid for which the volume $s$ remains unchanged, that is, for which
\[
\delta s=\int\mathrm{d}U.\delta e.\cos(4,5)
\]
vanishes. It is easy to see that in the state of equilibrium the shape of the surface $U$ must be such that at every point of it this element of the variation $\delta W$
\[
\mathrm{d}U.\delta e.\cos(4,5).\bigl[z+\alpha\alpha(\tfrac1R+\tfrac1{R'})\bigr]
\]
is proportional to the element of the variation $\delta s$, namely the quantity $\mathrm{d}U.\delta e.\cos(4,5)$, or equivalently that
\[
z+\alpha\alpha(\tfrac1R+\tfrac1{R'})=\mathrm{Const.}
\]
Clearly, if this proportionality did not hold, the value of $W$ could be decreased by a suitable change in the shape of the surface $U$, even while the boundary $P$ remained unchanged. Moreover, that equation holds for the \textit{whole} surface $U$, even if it consists of several separate parts, provided that the fluid itself is connected.

This equation constitutes the first fundamental theorem in the theory of the equilibrium of fluids. The illustrious \textsc{Laplace} had already obtained it by a method entirely different from ours.

If we take the plane for which $z$ equals the constant in the equation, which we may call the normal horizontal plane (plan de niveau), in place of the plane from which the coordinates $z$ were measured, we shall have
\[
z=-\alpha\alpha(\tfrac1R+\tfrac1{R'})
\]
which immediately yields the following corollaries.

I. If the normal plane cuts the free surface $U$ anywhere, the surface will necessarily be concave-convex at every point of the intersection, and the greatest radius of convexity will equal the greatest radius of concavity.

II. Above the normal plane the surface will either be concave-concave or, wherever it is concave-convex, the concave curvature will exceed the convex curvature.

III. Below the normal plane the surface will either be convex-convex or, wherever it is concave-convex, the convex curvature will exceed the concave curvature.

IV. The free surface $U$ can have a finite flat part only if that part is horizontal and coincides with the normal plane.

\SourceLeaf{73}\Article{29}
If the equation we have just established holds, the variation of the value of $W$ reduces to
\[
\delta W=-\int\mathrm{d}P.\delta e.\cos(5,8)(\alpha\alpha\cos i-\alpha\alpha+2\beta\beta)
\]
so that, on introducing an angle $A$ such that
\[
\cos A=\frac{\alpha\alpha-2\beta\beta}{\alpha\alpha}\qquad\text{or}\qquad\sin\tfrac12 A=\frac\beta\alpha
\]
we have
\[
\delta W=\alpha\alpha\int\mathrm{d}P.\delta e.\cos(5,8).(\cos A-\cos i)
\]
with the integration extending over the whole line $P$. We must remember that the factor $\cos(5,8)$ is equal to $\sin(5,6)$ with a positive or negative sign, according as we imagine the fluid in its virtual motion at the element $\mathrm{d}P$ to spill beyond the boundary $P$ or recede back within it. From this we readily conclude that in the state of equilibrium, generally speaking, $i=A$ must hold everywhere. For if at some part of the line $P$ we had $i<A$, a virtual motion of the first kind in that part, with the rest of the boundary $P$ remaining unchanged, would clearly give $W$ a negative variation. Likewise, a virtual motion of the second kind would yield a negative variation of $W$ if at any part of the line $P$ we had $i>A$. Both assumptions therefore conflict with the condition for a minimum in equilibrium.

This is the second fundamental theorem. It also enters the investigations of the illustrious \textsc{Laplace}, though we find no demonstration there from the principle of molecular forces.

